\section{Conclusiones y trabajo futuro}
\label{sec:conclusiones}

Hemos mostrado que una red convolucional pequeña, entrenada desde cero, reconoce un vocabulario
de ocho gestos más una clase nula con una exactitud del \AccTestFinal{} (IC del 95\,\%:
\AccTestFinalIC) en personas que nunca vio durante el entrenamiento. Esto supera con holgura el
objetivo del 85\,\%. Creemos que las claves del resultado son de diseño más que de capacidad:
clasificar
recortes de la mano en lugar del fotograma completo, construir la clase nula con manos de las
mismas imágenes y medir el rendimiento con una partición por persona.

El estudio de ablación, repetido con dos semillas, indica que la normalización por lotes es el
componente imprescindible: sin ella el entrenamiento es mucho más lento y, con el mismo número
de épocas, comete más del doble de errores (\PmcnemarSinBn). El \emph{dropout} y el aumento de
datos modifican la dinámica del entrenamiento como se espera de un regularizador, pero su efecto
sobre la exactitud final es del mismo orden que la variación entre ejecuciones idénticas.

Frente a un ResNet18 preentrenado, con \RatioParams{} veces más parámetros, la diferencia en test
no es estadísticamente significativa (\PmcnemarResnet). Ambos modelos clasifican un recorte
en pocos milisegundos incluso en CPU, de modo que cualquiera de los dos sirve para una aplicación
en tiempo real. Grad-CAM confirma que la red basa sus decisiones en la forma de la mano, y los
errores restantes se deben sobre todo a gestos realmente parecidos, a imágenes de baja calidad y a
ruido de etiqueta.

\paragraph{Trabajo futuro.}
\begin{itemize}
  \item \textbf{Medir el cambio de dominio.} Evaluar el sistema completo con las cajas de
    MediaPipe y con vídeos propios grabados con distintas iluminaciones, y comparar el resultado
    con el de test. Es el experimento que más falta para valorar el uso real.
  \item \textbf{Entrenar con cajas de detector.} Si el cambio de dominio es grande, entrenar con
    recortes generados por MediaPipe, o perturbar más las cajas durante el aumento de datos.
  \item \textbf{Más resolución.} Usar la versión de HaGRID a 512\,px o superior para conservar
    detalle en manos pequeñas, que es lo que distingue \cls{like} de \cls{call}.
  \item \textbf{Clase nula más rica.} Revisar las etiquetas ruidosas de \cls{no\_gesture} y
    añadir posturas de transición entre gestos, que son las que más riesgo tienen de activar una
    orden por error.
  \item \textbf{Robot real.} Pasar de \emph{turtlesim} a Gazebo o a un robot físico, que usa la
    misma interfaz de velocidades.
\end{itemize}
